\section{The actual native factor from any original model centre frame}

Fix the original test configuration and the same action/fibre-one data
on a line $M$ with original generic line $L$.  At the actual original
Gauss model point, let $e_C$ be any frame over its proper valuation centre.
Let $e_U$ be a frame of $L$ on an original open neighbourhood of the
original proper centre.  Define $c(e_C,e_U)$ by applying the actual
product residue valuation to the transition unit from the original
reference frame to the frame induced by $e_C$ and the given action data.

\begin{lemma}[Every actual centre frame computes the same section factor]
\label{lem:native-tc-nef-fs-arbitrary-center-factor-spec}
\leanok
\uses{thm:native-tc-coherent-all-power-section-norms}
\lean{EnvelopeProperty.NativeTestConfiguration.TestConfiguration.gaussPowerSectionNorm_eq_centerFactor}
For every degree $n$ and every original section $s$ of $M^{\otimes n}$,
its true original Gauss model norm equals
$c(e_C,e_U)^n$ times its original weighted Laurent norm.
\end{lemma}
\begin{proof}
The generic frame induced by a frame at the proper valuation centre,
and all its actual power frames, have true original model unit norm one.
The actual framed-pullback norm formula therefore computes the Gauss
model norm by the coefficient in that same model power frame.
The actual power-frame transition unit is the $n$th power of the
degree-one transition unit.  Changing to the original reference frame
multiplies coefficient values by this unit valuation.  Actual Laurent
support evaluation in the original proper-centre open identifies the
remaining coefficient value with the original weighted Laurent norm.
\end{proof}

\begin{lemma}[Independence from the original model centre frame]
\label{lem:native-tc-nef-fs-native-factor-arbitrary-center}
\leanok
\uses{lem:native-tc-nef-fs-arbitrary-center-factor-spec, thm:native-tc-unique-positive-all-power-factor, def:native-tc-intrinsic-native-unit-factor}
\lean{EnvelopeProperty.NativeTestConfiguration.TestConfiguration.gaussNativeNormFactor_eq_centerFactor}
If the original generic line $L$ is actually ample, the defined original
native norm factor equals $c(e_C,e_U)$ for every such actual model centre
frame and original reference frame.
\end{lemma}
\begin{proof}
Actual ampleness produces an original positive-degree section of positive
Gauss norm.  The original uniqueness theorem then compares the defined
native factor with the factor in the preceding lemma: both are positive
and satisfy the same equality for all original degrees and sections.
\end{proof}

\begin{remark}
The ample hypothesis is ordinary geometric ampleness of the original
generic line.  In the intended PL application it follows from the
positive original $\mathcal O(n)$ power.  No potential identity or
universal numerical-class invariance is assumed.
\end{remark}
